\chapter{LCS con al massimo 3 simboli rossi}

\section{Definizione del problema}

\subsection{Istanza}
Siano date due sequenze di elementi sull'alfabeto $\Sigma$:
\begin{align*}
    X &= \langle x_1, x_2, \dots, x_m \rangle \\
    Y &= \langle y_1, y_2, \dots, y_n \rangle
\end{align*}
e una funzione di colorazione $\text{col} : \Sigma \to \{\text{rosso}, \text{blu}, \text{nero}\}$ che associa a ciascun simbolo un colore.

\subsection{Soluzione}
Una più lunga sottosequenza comune tra $X$ e $Y$ la quale contiene al massimo $3$ elementi di colore rosso.

\section{Definizione del sottoproblema e del coefficiente}
Per ogni $i \in \{0, \dots, m\}$, per ogni $j \in \{0, \dots, n\}$ e per ogni $r \in \{0, 1, 2, 3\}$:
\begin{itemize}
    \item $X_i = \langle x_1, \dots, x_i \rangle$ è il prefisso di lunghezza $i$ di $X$;
    \item $Y_j = \langle y_1, \dots, y_j \rangle$ è il prefisso di lunghezza $j$ di $Y$;
    \item $r$ rappresenta il budget massimo di elementi rossi ammessi nella sottosequenza comune.
\end{itemize}

\medskip
\noindent
\textbf{Definizione del coefficiente $c_{i,j}^r$:} Per ogni $i \in \{0, \dots, m\}$, $j \in \{0, \dots, n\}$ e $r \in \{0, 1, 2, 3\}$, $c_{i,j}^r$ è la lunghezza di una più lunga sottosequenza comune tra i prefissi $X_i$ e $Y_j$ la quale contiene al massimo $r$ elementi di colore rosso.

\medskip
\noindent
\textbf{Numero totale di sottoproblemi:} $4 \cdot (m + 1) \cdot (n + 1) = \Theta(m \cdot n)$.

\section{Equazioni di ricorrenza}

\subsection{Caso base}
\[
    (i = 0 \lor j = 0) \land r \in \{0, 1, 2, 3\} \implies c_{i,j}^r = 0
\]
\textbf{Giustificazione:} Quando almeno uno dei due prefissi è vuoto, l'unica sottosequenza comune possibile è la sequenza vuota $\langle \rangle$, che contiene $0 \le r$ elementi rossi (rispettando il vincolo per ogni $r \ge 0$) ed ha lunghezza pari a $0$.

\subsection{Passo ricorsivo}
Per $i > 0 \land j > 0$ e per ogni $r \in \{0, 1, 2, 3\}$:
\begin{itemize}
    \item \textbf{Caso $x_i \neq y_j$:}
    \[
        c_{i,j}^r = \max \{ c_{i-1, j}^r, \, c_{i, j-1}^r \}
    \]
    \textbf{Giustificazione:} Poiché $x_i$ e $y_j$ sono diversi, non possono terminare simultaneamente la sottosequenza comune; si conserva il budget $r$ e si sceglie il massimo escludendo $x_i$ oppure $y_j$.

    \item \textbf{Caso $x_i = y_j$ con $\text{col}(x_i) \neq \text{rosso}$:}
    \[
        c_{i,j}^r = c_{i-1, j-1}^r + 1
    \]
    \textbf{Giustificazione:} Un match di colore blu o nero non consuma il budget dei simboli rossi. È quindi sempre ottimo includerlo nella soluzione aggiungendo $1$ alla lunghezza ottima del sottoproblema $(i-1, j-1, r)$.

    \item \textbf{Caso $x_i = y_j$ con $\text{col}(x_i) = \text{rosso}$:}
    \begin{itemize}
        \item Se $r = 0$:
        \[
            c_{i,j}^0 = c_{i-1, j-1}^0
        \]
        \textbf{Giustificazione:} Con budget $r = 0$ non è consentito prelevare alcun simbolo rosso; l'elemento $x_i = y_j$ viene obbligatoriamente scartato.
        \item Se $r > 0$:
        \[
            c_{i,j}^r = \max \{ c_{i-1, j-1}^{r-1} + 1, \, c_{i-1, j-1}^r \}
        \]
        \textbf{Giustificazione:} Avendo budget $r > 0$, si valuta il massimo tra due scelte: includere il simbolo rosso $x_i = y_j$ (aumentando di $1$ la lunghezza e riducendo il budget residuo a $r-1$) oppure scartarlo (mantenendo invariato il budget $r$ sul sottoproblema $(i-1, j-1, r)$).
    \end{itemize}
\end{itemize}

\subsection{Soluzione del problema principale}
La lunghezza della più lunga sottosequenza comune tra $X$ e $Y$ con al massimo $3$ elementi rossi è data da:
\[
    \text{Soluzione} = c_{m, n}^3
\]

\section{Algoritmo Bottom-Up}
\begin{algorithm}[htbp]
\caption{Calcolo Bottom-Up per LCS con al massimo 3 simboli rossi}
\begin{algorithmic}[1]
\Procedure{LCS3R}{$X, Y, \text{col}$}
    \State $m \gets |X|$
    \State $n \gets |Y|$
    \State $c \gets \text{nuova matrice } c[0 \dots m, \, 0 \dots n, \, 0 \dots 3]$
    \Comment{Caso base}
    \For{$i = 0 \textbf{ to } m$}
        \For{$r = 0 \textbf{ to } 3$}
            \State $c[i, 0, r] \gets 0$
        \EndFor
    \EndFor
    \For{$j = 0 \textbf{ to } n$}
        \For{$r = 0 \textbf{ to } 3$}
            \State $c[0, j, r] \gets 0$
        \EndFor
    \EndFor
    \Comment{Passo ricorsivo}
    \For{$i = 1 \textbf{ to } m$}
        \For{$j = 1 \textbf{ to } n$}
            \For{$r = 0 \textbf{ to } 3$}
                \If{$x_i \neq y_j$}
                    \State $c[i, j, r] \gets \max \{ c[i-1, j, r], \, c[i, j-1, r] \}$
                \Else
                    \If{$\text{col}(x_i) \neq \text{rosso}$}
                        \State $c[i, j, r] \gets c[i-1, j-1, r] + 1$
                    \Else
                        \If{$r == 0$}
                            \State $c[i, j, 0] \gets c[i-1, j-1, 0]$
                        \Else
                            \State $c[i, j, r] \gets \max \{ c[i-1, j-1, r-1] + 1, \, c[i-1, j-1, r] \}$
                        \EndIf
                    \EndIf
                \EndIf
            \EndFor
        \EndFor
    \EndFor
    \State \Return $c$
\EndProcedure
\end{algorithmic}
\end{algorithm}

\section{Algoritmo di ricostruzione}
\begin{algorithm}[htbp]
\caption{Ricostruzione della sequenza LCS con al massimo 3 rossi}
\begin{algorithmic}[1]
\Procedure{LCS3RR}{$c, X, Y, i, j, r$}
    \If{$i == 0 \lor j == 0$}
        \State \Return
    \Else
        \If{$x_i \neq y_j$}
            \If{$c[i-1, j, r] \ge c[i, j-1, r]$}
                \State \Call{LCS3RR}{$c, X, Y, i-1, j, r$}
            \Else
                \State \Call{LCS3RR}{$c, X, Y, i, j-1, r$}
            \EndIf
        \Else
            \If{$\text{col}(x_i) \neq \text{rosso}$}
                \State \Call{LCS3RR}{$c, X, Y, i-1, j-1, r$}
                \State \textbf{print} $x_i$
            \Else
                \If{$r == 0$}
                    \State \Call{LCS3RR}{$c, X, Y, i-1, j-1, 0$}
                \Else
                    \If{$c[i-1, j-1, r-1] + 1 \ge c[i-1, j-1, r]$}
                        \State \Call{LCS3RR}{$c, X, Y, i-1, j-1, r-1$}
                        \State \textbf{print} $x_i$
                    \Else
                        \State \Call{LCS3RR}{$c, X, Y, i-1, j-1, r$}
                    \EndIf
                \EndIf
            \EndIf
        \EndIf
    \EndIf
\EndProcedure
\end{algorithmic}
\end{algorithm}

\noindent
\textbf{Chiamata principale (Driver):} Per stampare la soluzione ottima del problema principale si invoca $\Call{LCS3RR}{c, X, Y, m, n, 3}$.

\section{Analisi della complessità}
\begin{itemize}
    \item \textbf{Complessità spaziale:} La tabella tridimensionale $c$ contiene $4 \cdot (m+1) \cdot (n+1)$ celle, pertanto richiede spazio $\Theta(m \cdot n)$.
    \item \textbf{Complessità temporale (Bottom-Up):} Ogni cella viene calcolata in tempo $\Theta(1)$ attraverso tre cicli annidati (di cui quello su $r$ esegue $4$ iterazioni costanti). Il tempo totale è $\Theta(4 \cdot m \cdot n) = \Theta(m \cdot n)$.
    \item \textbf{Complessità temporale (Ricostruzione):} Ad ogni passo ricorsivo almeno uno tra $i$ e $j$ diminuisce di $1$, pertanto la ricostruzione termina in tempo $O(m + n)$.
\end{itemize}